\chapter{Explaining the $\mathcal{P}(\mathbb{N})$ partition}

	\noindent
	This is easy. We are going to split $\mathcal{P}(\mathbb{N})$ into three subsets, formed by subsets of $\mathbb{N}$.
	\\\\
	
	\noindent
	The first one, only contains $\emptyset$.
	\\\\
	
	\noindent
	The second one, will contains all possible subsets of $\mathbb{N}$ with finite cardinal.
	\\\\
	
	\noindent
	The third one, will contains all possible subsets of $\mathbb{N}$ with INFINITE cardinal.
	\\\\
	
	\noindent
	In the general scheme \footnote{Figure \ref{fig:gereral_scheme}}, we can see them under "P(N)": the light blue squares. In the leftest side of the scheme.
	\\\\
	
	
	
\newpage